\documentclass[10pt,a4paper]{amsart}
\usepackage{amsmath,amssymb,amsthm,graphicx}
% Retained minimal preamble for the Stacks Project chapter more-algebra.tex (More on Algebra).
% The chapter file pinned-original-source.tex carries only \input{preamble} on its first line, so no
% preamble text lives inside the chapter itself.  The Stacks Project uses one single shared preamble for
% every chapter; the snapshot stages it identically as sources/stacks-chow/preamble.tex and
% sources/stacks-etale-cohomology/preamble.tex.  The block below is copied mechanically, line for line,
% from sources/stacks-chow/preamble.tex lines 14-18 (the Xy-pic setup) and 154-217 (the native theorem
% declarations together with the mathematical macros).  Nothing was added, removed or redefined; the
% wrapper class and the amsmath/amssymb/amsthm/graphicx base packages are supplied by the assembler.
% Chapter: more-algebra.tex -- More on Algebra.
% For commutative diagrams we use Xy-pic
\usepackage[all]{xy}

% We use 2cell for 2-commutative diagrams.
\xyoption{2cell}

% Theorem environments.
%
\theoremstyle{plain}
\newtheorem{theorem}[subsection]{Theorem}
\newtheorem{proposition}[subsection]{Proposition}
\newtheorem{lemma}[subsection]{Lemma}

\theoremstyle{definition}
\newtheorem{definition}[subsection]{Definition}
\newtheorem{example}[subsection]{Example}
\newtheorem{exercise}[subsection]{Exercise}
\newtheorem{situation}[subsection]{Situation}

\theoremstyle{remark}
\newtheorem{remark}[subsection]{Remark}
\newtheorem{remarks}[subsection]{Remarks}

\numberwithin{equation}{subsection}

% Macros
%
\def\lim{\mathop{\mathrm{lim}}\nolimits}
\def\colim{\mathop{\mathrm{colim}}\nolimits}
\def\Spec{\mathop{\mathrm{Spec}}}
\def\Hom{\mathop{\mathrm{Hom}}\nolimits}
\def\Ext{\mathop{\mathrm{Ext}}\nolimits}
\def\SheafHom{\mathop{\mathcal{H}\!\mathit{om}}\nolimits}
\def\SheafExt{\mathop{\mathcal{E}\!\mathit{xt}}\nolimits}
\def\Sch{\mathit{Sch}}
\def\Mor{\mathop{\mathrm{Mor}}\nolimits}
\def\Ob{\mathop{\mathrm{Ob}}\nolimits}
\def\Sh{\mathop{\mathit{Sh}}\nolimits}
\def\NL{\mathop{N\!L}\nolimits}
\def\CH{\mathop{\mathrm{CH}}\nolimits}
\def\proetale{{pro\text{-}\acute{e}tale}}
\def\etale{{\acute{e}tale}}
\def\QCoh{\mathit{QCoh}}
\def\Ker{\mathop{\mathrm{Ker}}}
\def\Im{\mathop{\mathrm{Im}}}
\def\Coker{\mathop{\mathrm{Coker}}}
\def\Coim{\mathop{\mathrm{Coim}}}

% Boxtimes
%
\DeclareMathSymbol{\boxtimes}{\mathbin}{AMSa}{"02}

%
% Macros for moduli stacks/spaces
%
\def\QCohstack{\mathcal{QC}\!\mathit{oh}}
\def\Cohstack{\mathcal{C}\!\mathit{oh}}
\def\Spacesstack{\mathcal{S}\!\mathit{paces}}
\def\Quotfunctor{\mathrm{Quot}}
\def\Hilbfunctor{\mathrm{Hilb}}
\def\Curvesstack{\mathcal{C}\!\mathit{urves}}
\def\Polarizedstack{\mathcal{P}\!\mathit{olarized}}
\def\Complexesstack{\mathcal{C}\!\mathit{omplexes}}
% \Pic is the operator that assigns to X its picard group, usage \Pic(X)
% \Picardstack_{X/B} denotes the Picard stack of X over B
% \Picardfunctor_{X/B} denotes the Picard functor of X over B
\def\Pic{\mathop{\mathrm{Pic}}\nolimits}
\def\Picardstack{\mathcal{P}\!\mathit{ic}}
\def\Picardfunctor{\mathrm{Pic}}
\def\Deformationcategory{\mathcal{D}\!\mathit{ef}}

% --- appended 2026-10-12 (Chief): the project own definitions of the reference/slogan/history
% display environments, copied mechanically from the shared Stacks preamble lines 8-12
% (\usepackage{verbatim} plus the three \newenvironment lines, which the chapter preamble keeps
% verbatim from the project preamble). Nothing added, removed or redefined: these environments
% carry no mathematical content and the project itself typesets them as comments. Without these
% three lines a counted statement whose own text uses them fails with Environment slogan undefined.
\usepackage{verbatim}
\newenvironment{reference}{\comment}{\endcomment}
\newenvironment{slogan}{\comment}{\endcomment}
\newenvironment{history}{\comment}{\endcomment}
\usepackage{cleveref}
\crefname{originallemma}{original lemma}{original lemmas}
\Crefname{originallemma}{Original lemma}{Original lemmas}
\makeatletter
\newlabel{algebra-lemma-characterize-finite-presentation}{{00QO}{source}{}{sourcepointer-1}{}}
\newlabel{algebra-lemma-characterize-finite-presentation@cref}{{[originallemma][0][]00QO}{[1][0][]source}{}{}{}}
\makeatother
\crefname{originallemma}{original lemma}{original lemmas}
\Crefname{originallemma}{Original lemma}{Original lemmas}
\makeatletter
\newlabel{algebra-lemma-smooth-strong-lift}{{07K4}{source}{}{sourcepointer-2}{}}
\newlabel{algebra-lemma-smooth-strong-lift@cref}{{[originallemma][0][]07K4}{[1][0][]source}{}{}{}}
\makeatother
\crefname{originallemma}{original lemma}{original lemmas}
\Crefname{originallemma}{Original lemma}{Original lemmas}
\makeatletter
\newlabel{lemma-absolutely-flat-over-field}{{30504}{source}{}{sourcepointer-3}{}}
\newlabel{lemma-absolutely-flat-over-field@cref}{{[originallemma][0][]30504}{[1][0][]source}{}{}{}}
\makeatother
\crefname{originallemma}{original lemma}{original lemmas}
\Crefname{originallemma}{Original lemma}{Original lemmas}
\makeatletter
\newlabel{lemma-class-weakly-etale-over-field}{{30846}{source}{}{sourcepointer-4}{}}
\newlabel{lemma-class-weakly-etale-over-field@cref}{{[originallemma][0][]30846}{[1][0][]source}{}{}{}}
\makeatother
\crefname{originallemma}{original lemma}{original lemmas}
\Crefname{originallemma}{Original lemma}{Original lemmas}
\makeatletter
\newlabel{lemma-max-weakly-etale-subalgebra}{{30892}{source}{}{sourcepointer-5}{}}
\newlabel{lemma-max-weakly-etale-subalgebra@cref}{{[originallemma][0][]30892}{[1][0][]source}{}{}{}}
\makeatother
\crefname{originallemma}{original lemma}{original lemmas}
\Crefname{originallemma}{Original lemma}{Original lemmas}
\makeatletter
\newlabel{lemma-when-weakly-etale}{{30454}{source}{}{sourcepointer-6}{}}
\newlabel{lemma-when-weakly-etale@cref}{{[originallemma][0][]30454}{[1][0][]source}{}{}{}}
\makeatother
\title{Functoriality, colimit and reduction properties of the maximal weakly etale subalgebra}
\author{The Stacks Project Authors}
\date{stacks-project:more-algebra.tex:local-snapshot-20261012; GNU FDL 1.2 or later}
\begin{document}
\maketitle
\setcounter{subsection}{1}
\noindent\textit{Editorial note.} All delivered mathematical text is the author's own, copied from the pinned Stacks Project chapter file more-algebra.tex; no statement or proof step was rewritten, reordered or removed.\par\medskip
\noindent\textit{Editorial note.} The chapter file pinned-original-source.tex carries only \textbackslash{}input\{preamble\} on its first line, so no preamble text lives inside the chapter. The retained file sources/stacks-more-algebra/retained-preamble.tex is a mechanical, line-for-line copy of the project's single shared preamble as staged in the snapshot; nothing was added, removed or redefined and no package was added.\par\medskip
\noindent\textit{Editorial note.} The source URL is the source-level Stacks Project URL, as required when more than one card is produced from one source; the per-claim tag is carried by claim\_key.\par\medskip
\noindent\textit{Editorial note.} The counted statement span closes at the author's \textbackslash{}end\{lemma\} on source line 30930. Its last enumerated item is the author's own placeholder '\textbackslash{}item add more here.', which is delivered verbatim rather than removed and carries no mathematical condition. The ten references inside the counted proof span are recorded as non-load-bearing navigation pointers.\par\medskip
\noindent\textit{Allowed standard prerequisites (editorial).}\begin{itemize}
\item The construction and universal property of the maximal weakly etale subalgebra B\_max(A). invoked at statement lines 30914-30916 and used throughout the counted proof span. Is the defining property from which each listed assertion is deduced.
\item Permanence of weakly etale maps under composition, base change and filtered colimits. used throughout the counted proof span. Supplies the stability that lets the maximal subalgebra be identified after each operation.
\end{itemize}\medskip
\ifdefined\hypertarget\hypertarget{sourcepointer-1}{}\fi
\noindent\textit{Source reference.} Original lemma 00QO: focused-seeds/AGP/sources/stacks-snapshot-20261011/tags.txt, source lines 974--974. Cross-chapter navigation pointer into the Stacks Project; the standard fact it supplies is applied inside the counted proof span and the argument resting on it stays there.\par\medskip
\ifdefined\hypertarget\hypertarget{sourcepointer-2}{}\fi
\noindent\textit{Source reference.} Original lemma 07K4: focused-seeds/AGP/sources/stacks-snapshot-20261011/tags.txt, source lines 9644--9644. Cross-chapter navigation pointer into the Stacks Project; the standard fact it supplies is applied inside the counted proof span and the argument resting on it stays there.\par\medskip
\ifdefined\hypertarget\hypertarget{sourcepointer-3}{}\fi
\noindent\textit{Source reference.} Original lemma 30504: sources/stacks-more-algebra/pinned-original-source.tex, source lines 30504--30504. Same-chapter navigation pointer; the statement it names is standard for the section and the argument that consumes it stays inside the counted proof span.\par\medskip
\ifdefined\hypertarget\hypertarget{sourcepointer-4}{}\fi
\noindent\textit{Source reference.} Original lemma 30846: sources/stacks-more-algebra/pinned-original-source.tex, source lines 30846--30846. Same-chapter navigation pointer; the statement it names is standard for the section and the argument that consumes it stays inside the counted proof span.\par\medskip
\ifdefined\hypertarget\hypertarget{sourcepointer-5}{}\fi
\noindent\textit{Source reference.} Original lemma 30892: sources/stacks-more-algebra/pinned-original-source.tex, source lines 30892--30892. Same-chapter navigation pointer; the statement it names is standard for the section and the argument that consumes it stays inside the counted proof span.\par\medskip
\ifdefined\hypertarget\hypertarget{sourcepointer-6}{}\fi
\noindent\textit{Source reference.} Original lemma 30454: sources/stacks-more-algebra/pinned-original-source.tex, source lines 30454--30454. Same-chapter navigation pointer; the statement it names is standard for the section and the argument that consumes it stays inside the counted proof span.\par\medskip
\begin{lemma}
\label{lemma-properties-of-max-weakly-etale-subalgebra}
Let $K$ be a field. For a $K$-algebra $A$ denote $B_{max}(A)$
the maximal weakly \'etale $K$-subalgebra of $A$ as in
Lemma \ref{lemma-max-weakly-etale-subalgebra}. Then
\begin{enumerate}
\item any $K$-algebra map $A' \to A$ induces a $K$-algebra map
$B_{max}(A') \to B_{max}(A)$,
\item if $A' \subset A$, then $B_{max}(A') = B_{max}(A) \cap A'$,
\item if $A = \colim A_i$ is a filtered colimit, then
$B_{max}(A) = \colim B_{max}(A_i)$,
\item the map $B_{max}(A) \to B_{max}(A_{red})$ is an isomorphism,
\item $B_{max}(A_1 \times \ldots \times A_n) =
B_{max}(A_1) \times \ldots \times B_{max}(A_n)$,
\item if $A$ has no nontrivial idempotents, then $B_{max}(A)$ is a
field and a separable algebraic extension of $K$,
\item add more here.
\end{enumerate}
\end{lemma}

\begin{proof}
Proof of (1). This is true because the image of $B_{max}(A') \to A$
is weakly \'etale over $K$ by Lemma \ref{lemma-class-weakly-etale-over-field}.

\medskip\noindent
Proof of (2). By (1) we have $B_{max}(A') \subset B_{max}(A)$.
Conversely, $B_{max}(A) \cap A'$ is a weakly \'etale $K$-algebra
by Lemma \ref{lemma-class-weakly-etale-over-field} and hence
contained in $B_{max}(A')$.

\medskip\noindent
Proof of (3). By (1) there is a map $\colim B_{max}(A_i) \to A$
which is injective because the system is filtered and
$B_{max}(A_i) \subset A_i$. The colimit $\colim B_{max}(A_i)$
is weakly \'etale over $K$ by
Lemma \ref{lemma-when-weakly-etale}.
Hence we get an injective map $\colim B_{max}(A_i) \to B_{max}(A)$.
Suppose that $a \in B_{max}(A)$. Then $a$ generates a finitely presented
$K$-subalgebra $B \subset B_{max}(A)$.
By Algebra, Lemma \ref{algebra-lemma-characterize-finite-presentation}
there is an $i$ and a $K$-algebra map $f : B \to A_i$
lifting the given map $B \to A$.
Since $B$ is weakly \'etale by
Lemma \ref{lemma-class-weakly-etale-over-field},
we see that $f(B) \subset B_{max}(A_i)$ and we conclude that $a$
is in the image of $\colim B_{max}(A_i) \to B_{max}(A)$.

\medskip\noindent
Proof of (4). Write $B_{max}(A_{red}) = \colim B_i$
as a filtered colimit of \'etale $K$-algebras
(Lemma \ref{lemma-absolutely-flat-over-field}).
By Algebra, Lemma \ref{algebra-lemma-smooth-strong-lift}
for each $i$ there is a $K$-algebra map $f_i : B_i \to A$
lifting the given map $B_i \to A_{red}$. It follows that the
canonical map $B_{max}(A_{red}) \to B_{max}(A)$ is surjective.
The kernel consists of nilpotent elements and hence is zero
as $B_{max}(A_{red})$ is reduced
(Lemma \ref{lemma-class-weakly-etale-over-field}).

\medskip\noindent
Proof of (5). Omitted.

\medskip\noindent
Proof of (6). Follows from Lemma \ref{lemma-class-weakly-etale-over-field}
part (4).
\end{proof}

\end{document}
